\section{Nuclear Medicine Cardiac Imaging}


\section*{Patient focus}

\subsubsection*{How cardiac nuclear results are judged}
For perfusion studies, stress adequacy, tracer, attenuation correction, and relative versus quantitative analysis affect interpretation. Viability and amyloid protocols answer different questions and have separate evidence.


This chapter explains the examination or clinical problem in accessible language. Patients should use it to prepare questions and understand the purpose, experience, benefits, and limitations of imaging; local instructions and the official report take priority.


\subsubsection*{What Is This Test?}
Nuclear cardiac stress tests use radioactive tracers (thallium, sestamibi) to show blood flow to the heart muscle at rest and stress. Areas with reduced flow during stress but normal flow at rest indicate coronary artery disease.

\subsubsection*{Alternative Names}
Nuclear stress test, myocardial perfusion imaging (MPI), cardiac PET, SPECT cardiac imaging, nuclear cardiology, sestamibi scan, thallium stress test


\subsubsection*{Principle}
Myocardial perfusion imaging uses Tc-99m sestamibi or tetrofosmin injected at stress and rest. The tracer distributes proportionally to myocardial blood flow. SPECT imaging detects perfusion defects: reversible defects indicate ischemia; fixed defects indicate infarct. ECG-gating evaluates wall motion and ejection fraction.

\subsubsection*{History}
Nuclear cardiology began with the use of thallium-201 for myocardial perfusion imaging in the 1970s. The development of technetium-99m labeled agents in the 1980s improved image quality and reduced radiation exposure.

\subsubsection*{Why This Test Is Ordered}
Your doctor orders a nuclear stress test (also called myocardial perfusion imaging) when you have chest pain, shortness of breath on exertion, or other symptoms suggesting your heart is not getting enough blood. It answers the question: is your coronary artery blood flow adequate at rest and during stress, and if not, which areas of your heart are affected?

\subsubsection*{Contraindications and Precautions}
Pregnancy and breastfeeding require examination-specific risk assessment, dose optimization, and current radiopharmaceutical guidance. For FDG studies, follow the protocol for preparation and glucose management.

\subsubsection*{Risks}
Risks associated with nuclear medicine include:
\begin{itemize}
    \item \textbf{Radiation exposure}: Low-dose radiation exposure from the administered radiopharmaceutical tracer.
    \item \textbf{Injection site discomfort}: Minor pain, bruising, or rare infection at the intravenous injection site.
\end{itemize}


You can resume normal activities. Results typically take one to three hours for immediate interpretation and one to three days for the formal report.

\subsubsection*{Understanding Your Results}

\subsubsection*{Normal Results}
A normal result shows uniform blood flow to the heart muscle at rest and during stress, with no reversible or fixed perfusion defects.

\begin{center}
\begin{tabular}{ll}
\toprule
\textbf{Population} & \textbf{Reference Range} \\
\midrule
Adult Female & No perfusion defects \\
Adult Male & No perfusion defects \\
Children & Apply age-specific imaging and safety criteria \\
Elderly & Apply age- and comorbidity-specific criteria \\
Pregnancy & Apply pregnancy-specific justification and safety criteria \\
\bottomrule
\end{tabular}
\end{center}
\subsubsection*{Follow-up Tests}
When a test result is not within the normal range, your doctor may recommend additional testing to confirm the finding, identify the underlying cause, and guide treatment decisions:
\begin{itemize}
    \item \textbf{Medication review and reconciliation:} Comprehensive review of all medications (prescribed, OTC, supplements) to identify drug interactions affecting laboratory results.
    \item \textbf{Serial monitoring:} Repeat testing at appropriate intervals to distinguish chronic stable abnormalities from acute changes requiring intervention.
    \item \textbf{Falls and functional assessment:} Evaluation of gait, balance, strength, and activities of daily living when abnormalities suggest functional decline.
    \item \textbf{Geriatric medicine consultation:} Comprehensive geriatric assessment for multi-morbidity management when abnormal results reflect complex age-related conditions.
\end{itemize}
\clearpage
